\documentclass[11pt
  , a4paper
  , article
  , oneside
  %  , twoside
  , showtrims
 % , draft
]{memoir}

\usepackage{essdocs}
\usepackage[numbers]{natbib}
\usepackage[autostyle]{csquotes}

\setsecnumdepth{subsection}

\begin{document}
%\frontmatter
%% ESS Document Description
%%
\essdocdesc{Engineering Manual}

%% ESS Document Number
%%
\essdocnum{ESS-XXXXXXXX}

%% Date
%%
\date{\today}

%% ESS Document Revision Number
%%
\essdocrev{0.1}

%% ESS Document State
%%
\essdocstate{Early Draft}

%% ESS Document Classification
%%
\essdocclass{Public}

%% Document Title
%%
\title{ICS Engineering Manual}
\subtitle{for an Inventory System with JIRA and EPICS}
%% Document Author(s), if more than one author,
%% use \newline instead of \\ or \linebreak in order to seperate them
\author{Jeong Han Lee (han.lee@esss.se)}

%% Document Reviewer(s) if more than one reviewer,
%% use \newline instead of \\ or \linebreak in order to seperate them
%\reviewer{Timo Korhonen (Chief Engineer) \newline Timo Korhonen (Chief Engineer)}
\reviewer{TBD}
%% Document Owner(s) if more than one owner,
%% use \newline instead of \\ or \linebreak in order to seperate them
\owner{ICS}

%% Document Approver(s) if more than one approver,
%% use \newline instead of \\ or \linebreak in order to seperate them
\approver{ICS}

\showtrimson

\esstitle
\newpage
\tableofcontents
\newpage

%\mainmatter


%%% Actual Document Start at below
\chapter{Overview}
There are infinitely ways to develop and maintain an inventory system. However, this inventory is the unique and \textbf{temporary} solution for the ICS Lab, especially, where the author works at. This inventory system is designed to minimize our time when we register an item to the existent JIRA ICS HW\&I group inventory task.  Therefore, it does \textbf{NOT} provide any fancy and beautiful ways to interact with users, BUT provide the minimal tool to monitor and track any equipment in ICS, ESS, and any IK partner. And the system only provide the following two features :
\begin{itemize}
\item Do stock an item to JIRA, and do Add-and-Print its barcodes.
\item Do add-and-print the barcodes of the item, which was registered in JIRA before.
\end{itemize}

\section{Inventory Workflow}

\subsection{How to stock an item at the first time to JIRA}

\subsection{How to add and print barcodes for an Item which has registered in JIRA}

\section{Components}
\subsection{Hardware}
\begin{itemize}
\item Honeywell Xenon 1900g (Wire) Barcode scanner or Honeywell Xenon 1902g (Wireless) Barcode scanner
\item DYMO LabelWriter 450 Duo
\end{itemize}
\subsection{Software}
\begin{itemize}
\item Linux OS
\item EPICS IOC
\item JIRA https://jira.esss.lu.se/projects/TAG/summary
\item DYMO LabelWriter 450 Duo CUP Driver
  http://www.dymo.com/en-US/dymo-label-sdk-and-cups-drivers-for-linux-dymo-label-sdk-cups-linux-p--1
\end{itemize}



    
\newpage

\chapter{Honeywell Xenon Scanners 1900 and 1902g}
Note that there are different procedure to setup Xenon 1900 corded scanner and Xenon 1902g cordless scanner. Please consult each manual in detail.


\section{Basic Configuration}
\begin{figure}[!htb]
%  \centering
  \includegraphics[width=0.3\textwidth]{./pictures/scanner_settings_default.eps}
  \caption{
            Default Settings.
          }
  \label{fig:default}   
\end{figure}

\begin{figure}[!htb]
%  \centering
  \includegraphics[width=0.3\textwidth]{./pictures/scanner_settings_usbserial.eps}
  \caption{
            USB Serial Setting.
          }
  \label{fig:usb-serial}   
\end{figure}


\begin{figure}[!htb]
%  \centering
  \includegraphics[width=0.3\textwidth]{./pictures/scanner_settings_silent_corded.eps}
  \caption{
            Silent Mode for Corded Scanner.
          }
  \label{fig:silent-corded}   
\end{figure}



\begin{figure}[!htb]
%  \centering
  \includegraphics[width=0.3\textwidth]{./pictures/scanner_settings_silent_cordless.eps}
  \caption{
            Silent Mode for Cordless Scanner.
          }
  \label{fig:silent-cordless}   
\end{figure}


\newpage

%% \chapter{Frequently Usage Case}
%% \section{MTCA}
%% \subsection{MTCA IOxOS IFC1410}
%% \subsection{MRF EVR-300DC}


%% \newpage
%% \section{Network Device}
%% \subsection{MOXA Nport 6650}
%% \subsection{HP Network Switch}

\newpage
\chapter{Predefined Bar Codes}
\section{Vendor Codes}
\vspace{1cm}
\noindent
\vspace{1.4cm}
\begin{minipage}{.2\textwidth}
\begin{center}
\includegraphics[height=1cm,keepaspectratio]{/home/jhlee/epics_env/epics-Apps/hw-xenon1900/docs/zint/vd_images/undefined.eps}
\end{center}
\end{minipage}
\begin{minipage}{.7\textwidth}
undefined
\end{minipage}


\noindent
\vspace{1.4cm}
\begin{minipage}{.2\textwidth}
\begin{center}
\includegraphics[height=1cm,keepaspectratio]{/home/jhlee/epics_env/epics-Apps/hw-xenon1900/docs/zint/vd_images/ess.eps}
\end{center}
\end{minipage}
\begin{minipage}{.7\textwidth}
ess
\end{minipage}


\noindent
\vspace{1.4cm}
\begin{minipage}{.2\textwidth}
\begin{center}
\includegraphics[height=1cm,keepaspectratio]{/home/jhlee/epics_env/epics-Apps/hw-xenon1900/docs/zint/vd_images/mrf.eps}
\end{center}
\end{minipage}
\begin{minipage}{.7\textwidth}
mrf
\end{minipage}


\noindent
\vspace{1.4cm}
\begin{minipage}{.2\textwidth}
\begin{center}
\includegraphics[height=1cm,keepaspectratio]{/home/jhlee/epics_env/epics-Apps/hw-xenon1900/docs/zint/vd_images/ioxos.eps}
\end{center}
\end{minipage}
\begin{minipage}{.7\textwidth}
ioxos
\end{minipage}


\noindent
\vspace{1.4cm}
\begin{minipage}{.2\textwidth}
\begin{center}
\includegraphics[height=1cm,keepaspectratio]{/home/jhlee/epics_env/epics-Apps/hw-xenon1900/docs/zint/vd_images/wiener.eps}
\end{center}
\end{minipage}
\begin{minipage}{.7\textwidth}
wiener
\end{minipage}


\noindent
\vspace{1.4cm}
\begin{minipage}{.2\textwidth}
\begin{center}
\includegraphics[height=1cm,keepaspectratio]{/home/jhlee/epics_env/epics-Apps/hw-xenon1900/docs/zint/vd_images/moxa.eps}
\end{center}
\end{minipage}
\begin{minipage}{.7\textwidth}
moxa
\end{minipage}


\noindent
\vspace{1.4cm}
\begin{minipage}{.2\textwidth}
\begin{center}
\includegraphics[height=1cm,keepaspectratio]{/home/jhlee/epics_env/epics-Apps/hw-xenon1900/docs/zint/vd_images/nat.eps}
\end{center}
\end{minipage}
\begin{minipage}{.7\textwidth}
nat
\end{minipage}


\noindent
\vspace{1.4cm}
\begin{minipage}{.2\textwidth}
\begin{center}
\includegraphics[height=1cm,keepaspectratio]{/home/jhlee/epics_env/epics-Apps/hw-xenon1900/docs/zint/vd_images/concurrent.eps}
\end{center}
\end{minipage}
\begin{minipage}{.7\textwidth}
concurrent
\end{minipage}


\noindent
\vspace{1.4cm}
\begin{minipage}{.2\textwidth}
\begin{center}
\includegraphics[height=1cm,keepaspectratio]{/home/jhlee/epics_env/epics-Apps/hw-xenon1900/docs/zint/vd_images/schroff.eps}
\end{center}
\end{minipage}
\begin{minipage}{.7\textwidth}
schroff
\end{minipage}


\noindent
\vspace{1.4cm}
\begin{minipage}{.2\textwidth}
\begin{center}
\includegraphics[height=1cm,keepaspectratio]{/home/jhlee/epics_env/epics-Apps/hw-xenon1900/docs/zint/vd_images/struck.eps}
\end{center}
\end{minipage}
\begin{minipage}{.7\textwidth}
struck
\end{minipage}


\noindent
\vspace{1.4cm}
\begin{minipage}{.2\textwidth}
\begin{center}
\includegraphics[height=1cm,keepaspectratio]{/home/jhlee/epics_env/epics-Apps/hw-xenon1900/docs/zint/vd_images/dell.eps}
\end{center}
\end{minipage}
\begin{minipage}{.7\textwidth}
dell
\end{minipage}


\noindent
\vspace{1.4cm}
\begin{minipage}{.2\textwidth}
\begin{center}
\includegraphics[height=1cm,keepaspectratio]{/home/jhlee/epics_env/epics-Apps/hw-xenon1900/docs/zint/vd_images/samsung.eps}
\end{center}
\end{minipage}
\begin{minipage}{.7\textwidth}
samsung
\end{minipage}


\noindent
\vspace{1.4cm}
\begin{minipage}{.2\textwidth}
\begin{center}
\includegraphics[height=1cm,keepaspectratio]{/home/jhlee/epics_env/epics-Apps/hw-xenon1900/docs/zint/vd_images/hp.eps}
\end{center}
\end{minipage}
\begin{minipage}{.7\textwidth}
hp
\end{minipage}


\noindent
\vspace{1.4cm}
\begin{minipage}{.2\textwidth}
\begin{center}
\includegraphics[height=1cm,keepaspectratio]{/home/jhlee/epics_env/epics-Apps/hw-xenon1900/docs/zint/vd_images/ibm.eps}
\end{center}
\end{minipage}
\begin{minipage}{.7\textwidth}
ibm
\end{minipage}


\noindent
\vspace{1.4cm}
\begin{minipage}{.2\textwidth}
\begin{center}
\includegraphics[height=1cm,keepaspectratio]{/home/jhlee/epics_env/epics-Apps/hw-xenon1900/docs/zint/vd_images/caen.eps}
\end{center}
\end{minipage}
\begin{minipage}{.7\textwidth}
caen
\end{minipage}


\noindent
\vspace{1.4cm}
\begin{minipage}{.2\textwidth}
\begin{center}
\includegraphics[height=1cm,keepaspectratio]{/home/jhlee/epics_env/epics-Apps/hw-xenon1900/docs/zint/vd_images/raritan.eps}
\end{center}
\end{minipage}
\begin{minipage}{.7\textwidth}
raritan
\end{minipage}


\noindent
\vspace{1.4cm}
\begin{minipage}{.2\textwidth}
\begin{center}
\includegraphics[height=1cm,keepaspectratio]{/home/jhlee/epics_env/epics-Apps/hw-xenon1900/docs/zint/vd_images/dymo.eps}
\end{center}
\end{minipage}
\begin{minipage}{.7\textwidth}
dymo
\end{minipage}




\newpage
\section{FormFactor Codes}
\vspace{1cm}
\noindent
\vspace{1.4cm}
\begin{minipage}{.2\textwidth}
\begin{center}
\includegraphics[height=1cm,keepaspectratio]{/home/iocuser/ics_gitsrc/epics_env/epics-Apps/hw-xenon1900/docs/zint/ff_images/undefined.eps}
\end{center}
\end{minipage}
\begin{minipage}{.7\textwidth}
undefined
\end{minipage}


\noindent
\vspace{1.4cm}
\begin{minipage}{.2\textwidth}
\begin{center}
\includegraphics[height=1cm,keepaspectratio]{/home/iocuser/ics_gitsrc/epics_env/epics-Apps/hw-xenon1900/docs/zint/ff_images/MTCA.eps}
\end{center}
\end{minipage}
\begin{minipage}{.7\textwidth}
MTCA
\end{minipage}


\noindent
\vspace{1.4cm}
\begin{minipage}{.2\textwidth}
\begin{center}
\includegraphics[height=1cm,keepaspectratio]{/home/iocuser/ics_gitsrc/epics_env/epics-Apps/hw-xenon1900/docs/zint/ff_images/PCIE.eps}
\end{center}
\end{minipage}
\begin{minipage}{.7\textwidth}
PCIE
\end{minipage}


\noindent
\vspace{1.4cm}
\begin{minipage}{.2\textwidth}
\begin{center}
\includegraphics[height=1cm,keepaspectratio]{/home/iocuser/ics_gitsrc/epics_env/epics-Apps/hw-xenon1900/docs/zint/ff_images/VME.eps}
\end{center}
\end{minipage}
\begin{minipage}{.7\textwidth}
VME
\end{minipage}


\noindent
\vspace{1.4cm}
\begin{minipage}{.2\textwidth}
\begin{center}
\includegraphics[height=1cm,keepaspectratio]{/home/iocuser/ics_gitsrc/epics_env/epics-Apps/hw-xenon1900/docs/zint/ff_images/FMC.eps}
\end{center}
\end{minipage}
\begin{minipage}{.7\textwidth}
FMC
\end{minipage}




%% \newpage
%% \section{ICS Location Codes}
%% \vspace{1cm}
%% \noindent
\vspace{1.4cm}
\begin{minipage}{.2\textwidth}
\begin{center}
\includegraphics[height=1cm,keepaspectratio]{/home/jhlee/epics_env/epics-Apps/hw-xenon1900/docs/zint/lo_images/undefined.eps}
\end{center}
\end{minipage}
\begin{minipage}{.7\textwidth}
undefined
\end{minipage}


\noindent
\vspace{1.4cm}
\begin{minipage}{.2\textwidth}
\begin{center}
\includegraphics[height=1cm,keepaspectratio]{/home/jhlee/epics_env/epics-Apps/hw-xenon1900/docs/zint/lo_images/icslab.eps}
\end{center}
\end{minipage}
\begin{minipage}{.7\textwidth}
icslab
\end{minipage}


\noindent
\vspace{1.4cm}
\begin{minipage}{.2\textwidth}
\begin{center}
\includegraphics[height=1cm,keepaspectratio]{/home/jhlee/epics_env/epics-Apps/hw-xenon1900/docs/zint/lo_images/utgard.eps}
\end{center}
\end{minipage}
\begin{minipage}{.7\textwidth}
utgard
\end{minipage}


\noindent
\vspace{1.4cm}
\begin{minipage}{.2\textwidth}
\begin{center}
\includegraphics[height=1cm,keepaspectratio]{/home/jhlee/epics_env/epics-Apps/hw-xenon1900/docs/zint/lo_images/site.eps}
\end{center}
\end{minipage}
\begin{minipage}{.7\textwidth}
site
\end{minipage}


\noindent
\vspace{1.4cm}
\begin{minipage}{.2\textwidth}
\begin{center}
\includegraphics[height=1cm,keepaspectratio]{/home/jhlee/epics_env/epics-Apps/hw-xenon1900/docs/zint/lo_images/ess.eps}
\end{center}
\end{minipage}
\begin{minipage}{.7\textwidth}
ess
\end{minipage}


\noindent
\vspace{1.4cm}
\begin{minipage}{.2\textwidth}
\begin{center}
\includegraphics[height=1cm,keepaspectratio]{/home/jhlee/epics_env/epics-Apps/hw-xenon1900/docs/zint/lo_images/office.eps}
\end{center}
\end{minipage}
\begin{minipage}{.7\textwidth}
office
\end{minipage}




%% \newpage
%% \section{Status Codes}
%% \vspace{1cm}
%% \noindent
\vspace{1.4cm}
\begin{minipage}{.2\textwidth}
\begin{center}
\includegraphics[height=1cm,keepaspectratio]{/home/jhlee/gitsrc/hw-xenon1900/docs/zint/st_images/undefined.eps}
\end{center}
\end{minipage}
\begin{minipage}{.7\textwidth}
undefined
\end{minipage}


\noindent
\vspace{1.4cm}
\begin{minipage}{.2\textwidth}
\begin{center}
\includegraphics[height=1cm,keepaspectratio]{/home/jhlee/gitsrc/hw-xenon1900/docs/zint/st_images/stock.eps}
\end{center}
\end{minipage}
\begin{minipage}{.7\textwidth}
stock
\end{minipage}


\noindent
\vspace{1.4cm}
\begin{minipage}{.2\textwidth}
\begin{center}
\includegraphics[height=1cm,keepaspectratio]{/home/jhlee/gitsrc/hw-xenon1900/docs/zint/st_images/inservice.eps}
\end{center}
\end{minipage}
\begin{minipage}{.7\textwidth}
inservice
\end{minipage}


\noindent
\vspace{1.4cm}
\begin{minipage}{.2\textwidth}
\begin{center}
\includegraphics[height=1cm,keepaspectratio]{/home/jhlee/gitsrc/hw-xenon1900/docs/zint/st_images/disposed.eps}
\end{center}
\end{minipage}
\begin{minipage}{.7\textwidth}
disposed
\end{minipage}





\newpage
\section{Model Codes}
\vspace{1cm}
\noindent
\vspace{1.4cm}
\begin{minipage}{.2\textwidth}
\begin{center}
\includegraphics[height=1cm,keepaspectratio]{/home/jhlee/epics_env/epics-Apps/hw-xenon1900/manual/zint/mo_images/undefined.eps}
\end{center}
\end{minipage}
\begin{minipage}{.7\textwidth}
UNDEFINED
\end{minipage}


\noindent
\vspace{1.4cm}
\begin{minipage}{.2\textwidth}
\begin{center}
\includegraphics[height=1cm,keepaspectratio]{/home/jhlee/epics_env/epics-Apps/hw-xenon1900/manual/zint/mo_images/mtca-evm-300.eps}
\end{center}
\end{minipage}
\begin{minipage}{.7\textwidth}
MTCA-EVM-300
\end{minipage}


\noindent
\vspace{1.4cm}
\begin{minipage}{.2\textwidth}
\begin{center}
\includegraphics[height=1cm,keepaspectratio]{/home/jhlee/epics_env/epics-Apps/hw-xenon1900/manual/zint/mo_images/vme-evm-300.eps}
\end{center}
\end{minipage}
\begin{minipage}{.7\textwidth}
VME-EVM-300
\end{minipage}


\noindent
\vspace{1.4cm}
\begin{minipage}{.2\textwidth}
\begin{center}
\includegraphics[height=1cm,keepaspectratio]{/home/jhlee/epics_env/epics-Apps/hw-xenon1900/manual/zint/mo_images/pcie-evr-300.eps}
\end{center}
\end{minipage}
\begin{minipage}{.7\textwidth}
PCIE-EVR-300
\end{minipage}


\noindent
\vspace{1.4cm}
\begin{minipage}{.2\textwidth}
\begin{center}
\includegraphics[height=1cm,keepaspectratio]{/home/jhlee/epics_env/epics-Apps/hw-xenon1900/manual/zint/mo_images/pcie-evr-300dc.eps}
\end{center}
\end{minipage}
\begin{minipage}{.7\textwidth}
PCIE-EVR-300DC
\end{minipage}


\noindent
\vspace{1.4cm}
\begin{minipage}{.2\textwidth}
\begin{center}
\includegraphics[height=1cm,keepaspectratio]{/home/jhlee/epics_env/epics-Apps/hw-xenon1900/manual/zint/mo_images/pmc-evr-230.eps}
\end{center}
\end{minipage}
\begin{minipage}{.7\textwidth}
PMC-EVR-230
\end{minipage}


\noindent
\vspace{1.4cm}
\begin{minipage}{.2\textwidth}
\begin{center}
\includegraphics[height=1cm,keepaspectratio]{/home/jhlee/epics_env/epics-Apps/hw-xenon1900/manual/zint/mo_images/vme-evr-230.eps}
\end{center}
\end{minipage}
\begin{minipage}{.7\textwidth}
VME-EVR-230
\end{minipage}


\noindent
\vspace{1.4cm}
\begin{minipage}{.2\textwidth}
\begin{center}
\includegraphics[height=1cm,keepaspectratio]{/home/jhlee/epics_env/epics-Apps/hw-xenon1900/manual/zint/mo_images/vme-evr-230rf.eps}
\end{center}
\end{minipage}
\begin{minipage}{.7\textwidth}
VME-EVR-230RF
\end{minipage}


\noindent
\vspace{1.4cm}
\begin{minipage}{.2\textwidth}
\begin{center}
\includegraphics[height=1cm,keepaspectratio]{/home/jhlee/epics_env/epics-Apps/hw-xenon1900/manual/zint/mo_images/cpci-evr-230.eps}
\end{center}
\end{minipage}
\begin{minipage}{.7\textwidth}
CPCI-EVR-230
\end{minipage}


\noindent
\vspace{1.4cm}
\begin{minipage}{.2\textwidth}
\begin{center}
\includegraphics[height=1cm,keepaspectratio]{/home/jhlee/epics_env/epics-Apps/hw-xenon1900/manual/zint/mo_images/vme-fout-12.eps}
\end{center}
\end{minipage}
\begin{minipage}{.7\textwidth}
VME-FOUT-12
\end{minipage}


\noindent
\vspace{1.4cm}
\begin{minipage}{.2\textwidth}
\begin{center}
\includegraphics[height=1cm,keepaspectratio]{/home/jhlee/epics_env/epics-Apps/hw-xenon1900/manual/zint/mo_images/cpci-fout-12.eps}
\end{center}
\end{minipage}
\begin{minipage}{.7\textwidth}
CPCI-FOUT-12
\end{minipage}


\noindent
\vspace{1.4cm}
\begin{minipage}{.2\textwidth}
\begin{center}
\includegraphics[height=1cm,keepaspectratio]{/home/jhlee/epics_env/epics-Apps/hw-xenon1900/manual/zint/mo_images/cpci-evg-220.eps}
\end{center}
\end{minipage}
\begin{minipage}{.7\textwidth}
CPCI-EVG-220
\end{minipage}


\noindent
\vspace{1.4cm}
\begin{minipage}{.2\textwidth}
\begin{center}
\includegraphics[height=1cm,keepaspectratio]{/home/jhlee/epics_env/epics-Apps/hw-xenon1900/manual/zint/mo_images/cpci-evg-230.eps}
\end{center}
\end{minipage}
\begin{minipage}{.7\textwidth}
CPCI-EVG-230
\end{minipage}


\noindent
\vspace{1.4cm}
\begin{minipage}{.2\textwidth}
\begin{center}
\includegraphics[height=1cm,keepaspectratio]{/home/jhlee/epics_env/epics-Apps/hw-xenon1900/manual/zint/mo_images/vme-evg-230.eps}
\end{center}
\end{minipage}
\begin{minipage}{.7\textwidth}
VME-EVG-230
\end{minipage}


\noindent
\vspace{1.4cm}
\begin{minipage}{.2\textwidth}
\begin{center}
\includegraphics[height=1cm,keepaspectratio]{/home/jhlee/epics_env/epics-Apps/hw-xenon1900/manual/zint/mo_images/nport5510.eps}
\end{center}
\end{minipage}
\begin{minipage}{.7\textwidth}
NPORT5510
\end{minipage}


\noindent
\vspace{1.4cm}
\begin{minipage}{.2\textwidth}
\begin{center}
\includegraphics[height=1cm,keepaspectratio]{/home/jhlee/epics_env/epics-Apps/hw-xenon1900/manual/zint/mo_images/nport6650.eps}
\end{center}
\end{minipage}
\begin{minipage}{.7\textwidth}
NPORT6650
\end{minipage}


\noindent
\vspace{1.4cm}
\begin{minipage}{.2\textwidth}
\begin{center}
\includegraphics[height=1cm,keepaspectratio]{/home/jhlee/epics_env/epics-Apps/hw-xenon1900/manual/zint/mo_images/wienercrate450.eps}
\end{center}
\end{minipage}
\begin{minipage}{.7\textwidth}
WIENERCRATE450
\end{minipage}


\noindent
\vspace{1.4cm}
\begin{minipage}{.2\textwidth}
\begin{center}
\includegraphics[height=1cm,keepaspectratio]{/home/jhlee/epics_env/epics-Apps/hw-xenon1900/manual/zint/mo_images/pdu-x3-3600.eps}
\end{center}
\end{minipage}
\begin{minipage}{.7\textwidth}
PDU-X3-3600
\end{minipage}





\newpage
\chapter{Action Bar Codes}
\vspace{1cm}
\noindent
\vspace{1.4cm}
\begin{minipage}{.2\textwidth}
\begin{center}
\includegraphics[height=1cm,keepaspectratio]{/home/jhlee/epics_env/epics-Apps/hw-xenon1900/docs/zint/act_images/CL.eps}
\end{center}
\end{minipage}
\begin{minipage}{.7\textwidth}
Clear all scanned PVs
\end{minipage}


\noindent
\vspace{1.4cm}
\begin{minipage}{.2\textwidth}
\begin{center}
\includegraphics[height=1cm,keepaspectratio]{/home/jhlee/epics_env/epics-Apps/hw-xenon1900/docs/zint/act_images/00.eps}
\end{center}
\end{minipage}
\begin{minipage}{.7\textwidth}
----------- place holder ----------
\end{minipage}


\noindent
\vspace{1.4cm}
\begin{minipage}{.2\textwidth}
\begin{center}
\includegraphics[height=1cm,keepaspectratio]{/home/jhlee/epics_env/epics-Apps/hw-xenon1900/docs/zint/act_images/LE.eps}
\end{center}
\end{minipage}
\begin{minipage}{.7\textwidth}
Enable Label Printing after JIRA action
\end{minipage}


\noindent
\vspace{1.4cm}
\begin{minipage}{.2\textwidth}
\begin{center}
\includegraphics[height=1cm,keepaspectratio]{/home/jhlee/epics_env/epics-Apps/hw-xenon1900/docs/zint/act_images/LD.eps}
\end{center}
\end{minipage}
\begin{minipage}{.7\textwidth}
Disable Label Printing after JIRA action
\end{minipage}


\noindent
\vspace{1.4cm}
\begin{minipage}{.2\textwidth}
\begin{center}
\includegraphics[height=1cm,keepaspectratio]{/home/jhlee/epics_env/epics-Apps/hw-xenon1900/docs/zint/act_images/LP.eps}
\end{center}
\end{minipage}
\begin{minipage}{.7\textwidth}
Print the last existent label
\end{minipage}


\noindent
\vspace{1.4cm}
\begin{minipage}{.2\textwidth}
\begin{center}
\includegraphics[height=1cm,keepaspectratio]{/home/jhlee/epics_env/epics-Apps/hw-xenon1900/docs/zint/act_images/00.eps}
\end{center}
\end{minipage}
\begin{minipage}{.7\textwidth}
----------- place holder ----------
\end{minipage}


\noindent
\vspace{1.4cm}
\begin{minipage}{.2\textwidth}
\begin{center}
\includegraphics[height=1cm,keepaspectratio]{/home/jhlee/epics_env/epics-Apps/hw-xenon1900/docs/zint/act_images/JC.eps}
\end{center}
\end{minipage}
\begin{minipage}{.7\textwidth}
Create an JIRA issue 
\end{minipage}


\noindent
\vspace{1.4cm}
\begin{minipage}{.2\textwidth}
\begin{center}
\includegraphics[height=1cm,keepaspectratio]{/home/jhlee/epics_env/epics-Apps/hw-xenon1900/docs/zint/act_images/JU.eps}
\end{center}
\end{minipage}
\begin{minipage}{.7\textwidth}
Update an JIRA issue (Define the issue number first)
\end{minipage}


\noindent
\vspace{1.4cm}
\begin{minipage}{.2\textwidth}
\begin{center}
\includegraphics[height=1cm,keepaspectratio]{/home/jhlee/epics_env/epics-Apps/hw-xenon1900/docs/zint/act_images/JD.eps}
\end{center}
\end{minipage}
\begin{minipage}{.7\textwidth}
Delete an JIRA issue (Define the issue number first)
\end{minipage}





\chapter{Outlook}

\end{document}
